\section{摄政之后，西北成为一张持续的账单}

1022年真宗去世，年幼的赵祯即位，刘太后临朝。继位和摄政可由《宋史》仁宗本纪、章献明肃刘皇后传与《续资治通鉴长编》乾兴元年条互证；不过，把近十年政务逐项归为太后个人意志，正是后世宫廷叙事容易造成的错觉。她所维持的是一个须同时支付河北防务、辽宋盟约和官僚日常运转的中枢。1033年刘太后去世，仁宗亲政，改变的是决断与人事的重心，并不是一夜之间出现了不受辅臣、台谏和内廷约束的皇帝。

西北的变化使这种继承很快受到检验。1034年李元昊改国号、年号，1038年称帝，意味着党项首领不再满足于宋朝册封秩序。《宋史》夏国传、《长编》与《宋会要辑稿·蕃夷》能确定宋廷所面对的外交转折；它们却主要记录宋朝如何命名对手。西夏的建国礼、内部接受和财政基础，不能只从敌国战报反推，仍须由黑水城文书、碑刻和区域考古补足。对开封而言，关键并非一句称号是否“僭”，而是河套、河西道路、堡寨与水源附近，出现了一个能独立调动资源的政权。

1040年的三川口、1041年的好水川和1042年的定川寨，连续打破宋军速胜的期待。宋方本纪、范雍传、任福传和《长编》保留了败讯，却把将领责任、诱敌经过与伤亡数字包裹在战后追责中；不宜把奏报的军数当作已被核实的战场规模。三场战争共同呈现的约束较为清楚：黄土高原上的行军、给水、粮道和情报，限制了从东京发出的命令；西夏军队则能够把机动和地形结合为防御。于是边防不只是地图上一条线，而是一串需要驻兵、修寨、运粮和轮换的支点。

\section{改革的尺度与边境的尺度}

这种压力解释了1043年庆历新政何以在西夏战争中提出。范仲淹、富弼等人针对选任、考课与学校提出措施，希望重整官僚的入口和责任；《宋会要辑稿·职官》庆历条与《长编》可以追索政策的提出。它们证明诏令与议论，却不足以证明州县已经同样执行。到1045年，主要倡议者相继离开中枢，改革的整套推进停止。若把这段经历简化为某一派“胜”或“败”，便看不见考课、荐举和既得利益如何在不同层次制造阻力；它留下的是一批未消失的问题和一份日后反复被援引的改革记忆。

1044年宋夏和议则表明，战争与改革都不能脱离财政的限度。双方以赐予、互市和停战安排降低边境消耗。《宋史》仁宗本纪、夏国传和会要的和议条足以确认框架；“赐予”这一宋方用语不能自动换算为单向臣属，条款、称谓和实际执行应分开处理。和平没有取消堡寨、侦察和贸易摩擦，却让朝廷可以重新配置西北支出。这是以军备支撑的停战，而不是没有代价的安宁。

边疆也不只在西北。1052年侬智高在岭南与交趾边境起兵，活动及邕州、广州附近；次年狄青率军平定。《宋史》侬智高传和狄青传提供宋方的行动线索，但其族群标签和“平定”后的范围，都不等于当地社会的自我分类和全部恢复。广源一带的首领原在宋、交趾与山地网络之间争取资源与承认，朝廷调兵虽重获主要城镇，却没有消除跨境道路、地方武装和交换关系。仁宗朝的边防由此显出双重面貌：中枢可以集中军费和将领，却不能把边地的多重联系压缩成一纸州县名册。
